%============================================================
% CH1 · PATTERN 1 — SINGLE-PASS AGGREGATES: THEORY
%============================================================
\subsection{Pattern Theory}
\label{ssec:ch1-p1-theory}

\begin{patternbox}{Pattern 1 at a Glance}
Many array questions ask for a single number --- or one number per position --- that depends on \textbf{the multiset of values}, not on how the elements are wired together: the maximum, the second maximum, how often something occurs, the best pair ``earlier $\le$ later''. Then we read the array \textbf{once}, keep a few variables that summarise the prefix already read, and update them with each new element.

\begin{center}
\begin{tikzpicture}
    \dparray{g}{0}{}{4/adone, 9/adone, 4/adone, 7/adone, 1/aread, 9/acell, 6/acell}
    \aptrup{g4}{x = a[i]}{accentorange}
    \arange{g0}{g3}{summarised prefix $[0, i)$}{tabulationblue}
    \node[rectangle, rounded corners=3pt, draw=successgreen, fill=solutionbg, font=\scriptsize, align=left,
          anchor=west, text width=3.9cm] (st) at (7.0,0)
        {\textbf{state} (three largest distinct)\\ $m_1 = 9,\ m_2 = 7,\ m_3 = 4$\\[2pt]
         $x = 1$: smaller than $m_3$ --- no change};
    \draw[dep] (g4.east) to[bend right=25] (st.west);
\end{tikzpicture}
\end{center}

\begin{recurrencebox}
\centering
\textit{state} $\gets$ summary of the empty prefix\\[2pt]
for $i = 0, 1, \dots, n-1$:\quad \textit{state} $\gets$ \textsc{update}(\textit{state}, $a_i$)\\[2pt]
answer $=$ \textsc{read}(\textit{state}) \qquad $\bigO{n}$ time, $\bigO{1}$ extra memory
\end{recurrencebox}
Typical states: a sum or a count; the minimum or maximum with its index; the $k$ largest (distinct or not) for a small $k$; the best value of a pair formed with an earlier element; the length of the current streak. Sometimes a second pass reuses the state at every position (Problem~\ref{ssec:ch1-advantage}).
\end{patternbox}

%------------------------------------------------------------
\subsubsection*{From Part A to Pattern 1}

Lesson~\ref{ssec:ch1-scanning} introduced the scan with running state and its invariant ``the state summarises $[0, i)$''. This pattern is the \emph{catalogue}: which summaries can be updated in $\bigO{1}$ per element, how to initialise them without breaking on extreme values, and how to read the answer. It is the first thing to try on any array question whose answer is one number --- before sorting ($\bigO{n \log n}$), before hashing, before anything with two indices.

\begin{diagrambox}
\begin{tikzpicture}[x=1cm, y=1cm,
    q/.style={rectangle, rounded corners=3pt, draw=deepblue, fill=theorybg, line width=0.9pt,
              text width=2.6cm, align=center, minimum height=13mm, font=\scriptsize}]
    \node[q] (s) at (0,0) {\textbf{1. STATE}\\ what do I need to know about $[0, i)$?};
    \node[q] (i) at (3.2,0) {\textbf{2. INITIAL}\\ the summary of the empty prefix};
    \node[q] (u) at (6.4,0) {\textbf{3. UPDATE}\\ old state $+\ a_i \to$ new state, in $\bigO{1}$};
    \node[q] (r) at (9.6,0) {\textbf{4. READ}\\ answer from the final state (or during the scan)};
    \foreach \x/\y in {s/i, i/u, u/r} \draw[flow] (\x) -- (\y);
    \node[note, anchor=north, text width=11cm] at (4.8,-0.95) {if step 3 is impossible in $\bigO{1}$ (the median, ``the most frequent value'' with large values), the problem is not Pattern 1: sort, hash or count by value (Pattern~5)};
\end{tikzpicture}
\figcap{The four decisions of every single-pass aggregate. Steps 2 and 3 are where the bugs live.}
\end{diagrambox}

%------------------------------------------------------------
\subsubsection*{When to recognise it --- and when not}

Use Pattern 1 when
\begin{itemize}[leftmargin=*]
    \item the answer is a \textbf{maximum, minimum, count, sum} or ``the $k$-th largest'' for a \textbf{small constant} $k$ (Problem~\ref{ssec:ch1-third-maximum});
    \item each position needs a \textbf{global} value with itself excluded (``the strongest \emph{other} participant'', ``the sum of the others''): one pass for the global summary, a second pass to answer (Problem~\ref{ssec:ch1-advantage});
    \item the answer is the best \textbf{pair $i < j$} and the best partner of $j$ is an extreme of the prefix (``largest $a_j - a_i$'': the partner is the minimum so far; Exercise~E1.3);
    \item the answer is about \textbf{streaks}: the longest block of equal or consecutive values, the number of ``hills'' (Exercises E1.1, E1.6);
    \item the statement says ``$\bigO{n}$ time'' and ``$\bigO{1}$ extra memory''.
\end{itemize}

\textbf{Do NOT use it} when the update needs more than a constant amount of information about the prefix:
\begin{itemize}[leftmargin=*]
    \item the \textbf{median}, or the $k$-th largest for a $k$ that grows with $n$ --- sort (Chapter~\ref{ch:sorting}) or use a heap (the \textit{Stacks, Queues \& Heaps} book);
    \item ``\textbf{how many times} does each value occur'' --- a frequency array (Pattern~5) or a hash map (Chapter~\ref{ch:hashing});
    \item \textbf{subarray sums} for many ranges --- prefix sums (Chapter~\ref{ch:prefix_sums}); the maximum subarray sum is Kadane's algorithm, a DP (\dpbook{1}{Pattern 5}).
\end{itemize}

\begin{diagrambox}
\begin{tikzpicture}
    \begin{scope}
        \dparray{m}{0}{prefix}{1/adone, 5/adone, 3/adone}
        \node[note, anchor=west, text width=2.3cm, align=left] at (2.9,0) {median $3$};
    \end{scope}
    \begin{scope}[xshift=6.3cm]
        \dparray{q}{0}{}{1/adone, 5/adone, 3/adone, 9/aread, 8/aread}
        \node[note, anchor=north, text width=4.4cm] at (2.25,-0.55) {append $9, 8$: new median $5$ --- an element that was \emph{not} the old median};
    \end{scope}
    \node[note, text width=11cm] at (5.4,-1.95) {the summary ``median $= 3$'' does not tell us that a $5$ was seen; a correct update would need the whole sorted prefix};
\end{tikzpicture}
\figcap{Why the median is not a single-pass aggregate: its update needs more than the old median and the new element.}
\end{diagrambox}

%------------------------------------------------------------
\subsubsection*{The framework and why it is correct}

\begin{table}[H]
\centering\small
\renewcommand{\arraystretch}{1.25}
\begin{tabularx}{\linewidth}{@{}>{\raggedright\arraybackslash}p{3.1cm} >{\raggedright\arraybackslash}p{2.6cm} >{\raggedright\arraybackslash}X >{\raggedright\arraybackslash}p{2.4cm}@{}}
\toprule
\textbf{aggregate} & \textbf{state} & \textbf{update with $x = a_i$} & \textbf{empty prefix} \\
\midrule
sum, count & $s$ & $s \gets s + x$ (or $+[\,x \text{ qualifies}\,]$) & $0$ \\
maximum + first index & $m, p$ & if $x > m$: $m \gets x,\ p \gets i$ & $m = -\infty$ \\
maximum + last index & $m, p$ & if $x \ge m$: $m \gets x,\ p \gets i$ & $m = -\infty$ \\
max and how often & $m, c$ & $x > m$: $(x, 1)$;\ $x = m$: $c{+}{+}$ & $(-\infty, 0)$ \\
two largest (duplicates allowed) & $m_1 \ge m_2$ & $x > m_1$: shift $m_2 \gets m_1,\ m_1 \gets x$; else $x > m_2$: $m_2 \gets x$ & $-\infty, -\infty$ \\
$k$ largest \emph{distinct} & $m_1 > \dots > m_k$ & skip $x$ if equal to some $m_j$; else insert $x$ and drop the last & all $-\infty$ \\
best pair $a_j - a_i$, $i < j$ & $\mathit{lo}$, \textit{best} & $\mathit{best} \gets \max(\mathit{best}, x - \mathit{lo})$, \emph{then} $\mathit{lo} \gets \min(\mathit{lo}, x)$ & $\mathit{lo} = +\infty$ \\
longest streak & $\mathit{run}$, \textit{best} & $\mathit{run} \gets (x \text{ continues}) \;?\; \mathit{run}+1 : 1$; $\mathit{best} \gets \max(\mathit{best}, \mathit{run})$ & $0, 0$ \\
\bottomrule
\end{tabularx}
\end{table}

\begin{invariantbox}[Invariant of a single-pass aggregate]
Let $f$ be the quantity we want, defined on every prefix. Suppose there is an \textsc{update} with
\[
f(a_0, \dots, a_{i}) \;=\; \textsc{update}\big(f(a_0, \dots, a_{i-1}),\ a_i\big) \qquad\text{for every } i,
\]
and \textit{state} starts as $f(\text{empty prefix})$. \textbf{Invariant:} before step $i$, $\textit{state} = f(a_0, \dots, a_{i-1})$.

\textbf{Proof.} Induction on $i$. It holds for $i = 0$ by the choice of the initial value. If it holds before step $i$, the step sets $\textit{state} \gets \textsc{update}(f(a_0..a_{i-1}), a_i) = f(a_0..a_i)$, so it holds before step $i+1$. After step $n-1$, $\textit{state} = f(a_0, \dots, a_{n-1})$. \hfill$\blacksquare$

\textbf{Proving an update.} For each row of the table, check the equation by cases. Example, two largest values with duplicates allowed: the two largest of $P \cup \{x\}$ are the two largest among $\{m_1, m_2, x\}$, because no other element of $P$ can beat $m_2$. Comparing $x$ with $m_1$ and $m_2$ sorts these three values, which is exactly the shift rule.
\end{invariantbox}

%------------------------------------------------------------
\subsubsection*{The $k$ largest: shifting a small sorted list}

The ``two largest'' and ``three largest'' states are a sorted list of length $k$ kept by \emph{insertion}: find where $x$ belongs, shift the smaller ones down by one, and let the last fall off. For $k = 3$ there are four cases:

\begin{diagrambox}
\begin{tikzpicture}[x=1cm, y=1cm,
    sl/.style={cell, minimum width=9mm},
    sn/.style={cell, minimum width=9mm, fill=currentfill, draw=accentorange, line width=1pt},
    so/.style={cell, minimum width=9mm, fill=lightgray, text=midgray}]
    \foreach \x/\lab [count=\k from 0] in {12/{$x > m_1$}, 8/{$m_1 > x > m_2$}, 5/{$m_2 > x > m_3$}, 2/{$x < m_3$}}{
        \begin{scope}[xshift=\k*3.25cm]
            \node[font=\scriptsize\bfseries, text=deepblue] at (1.0,1.25) {\lab};
            \node[sl] at (0.1,0.55) {10}; \node[sl] at (1.0,0.55) {6}; \node[sl] at (1.9,0.55) {3};
            \node[font=\scriptsize, text=darkgray] at (1.0,0.0) {insert $x = \x$};
        \end{scope}
    }
    \begin{scope}[yshift=-0.55cm]
        \node[sn] at (0.1,0) {12}; \node[sl] at (1.0,0) {10}; \node[sl] at (1.9,0) {6};
        \begin{scope}[xshift=3.25cm] \node[sl] at (0.1,0) {10}; \node[sn] at (1.0,0) {8}; \node[sl] at (1.9,0) {6}; \end{scope}
        \begin{scope}[xshift=6.5cm] \node[sl] at (0.1,0) {10}; \node[sl] at (1.0,0) {6}; \node[sn] at (1.9,0) {5}; \end{scope}
        \begin{scope}[xshift=9.75cm] \node[sl] at (0.1,0) {10}; \node[sl] at (1.0,0) {6}; \node[sl] at (1.9,0) {3}; \end{scope}
    \end{scope}
\end{tikzpicture}
\figcap{Inserting $x$ into $(m_1, m_2, m_3) = (10, 6, 3)$. The value $3$ falls off in the first three cases; in the last one nothing changes.}
\end{diagrambox}

\textbf{Distinct or not?} Read the statement twice. ``The second largest \emph{element}'' of $(5, 5, 3)$ is $5$ (duplicates count separately); ``the second largest \emph{distinct value}'' is $3$. The update differs in one line: for distinct values, an $x$ equal to some $m_j$ is skipped \emph{before} the comparisons; for the non-distinct version, ties fall through to a lower slot.

\begin{diagrambox}
\begin{tikzpicture}
    \dparray{d}{0}{input}{5/acell, 5/acell, 3/acell}
    \begin{scope}[xshift=4.2cm]
        \node[font=\scriptsize\bfseries, text=deepblue, anchor=west] at (0,0.55) {duplicates allowed};
        \node[cell, fill=basefill] at (0.45,0) {5}; \node[cell, fill=basefill] at (1.35,0) {5};
        \node[font=\scriptsize, anchor=west] at (1.9,0) {2nd $= 5$};
    \end{scope}
    \begin{scope}[xshift=8.4cm]
        \node[font=\scriptsize\bfseries, text=deepblue, anchor=west] at (0,0.55) {distinct values};
        \node[cell, fill=basefill] at (0.45,0) {5}; \node[cell, fill=basefill] at (1.35,0) {3};
        \node[font=\scriptsize, anchor=west] at (1.9,0) {2nd $= 3$};
    \end{scope}
\end{tikzpicture}
\figcap{The same input, two meanings of ``second largest''. Problem~\ref{ssec:ch1-advantage} needs the left one, Problem~\ref{ssec:ch1-third-maximum} the right one.}
\end{diagrambox}

For a general small $k$ the insertion costs $\bigO{k}$ per element, so $\bigO{nk}$ in total --- still linear for $k = 3$. For large $k$ use a heap of size $k$ ($\bigO{n \log k}$) or \texttt{nth\_element} (average $\bigO{n}$).

%------------------------------------------------------------
\subsubsection*{Initial values and sentinels}

The empty prefix has no maximum. Three ways to represent ``nothing seen yet'', in order of preference:
\begin{enumerate}[leftmargin=*]
    \item \textbf{Start from the first element} when $n \ge 1$ is guaranteed: $m = a_0$, then loop from $i = 1$. No sentinel at all.
    \item \textbf{A sentinel outside the value range:} if values are \texttt{int}, keep the state in \texttt{long long} and start from \texttt{LLONG\_MIN}; no \texttt{int} can equal it. Starting from \texttt{INT\_MIN} is wrong when the data may contain \texttt{INT\_MIN} itself --- you cannot tell ``unset'' from ``set to $-2^{31}$'' (Problem~\ref{ssec:ch1-third-maximum} has exactly this trap).
    \item \textbf{An explicit flag or counter:} \texttt{bool have2} or ``number of slots filled''. Always correct, slightly more code.
\end{enumerate}
Starting a maximum from $0$ is the most common bug in this pattern: it is correct only when all values are known to be non-negative.

%------------------------------------------------------------
\subsubsection*{Pairs with an earlier element}

``Maximise $a_j - a_i$ over $i < j$'' looks like a double loop. Fix the \emph{later} index $j$: its best partner is the smallest value in $[0, j)$, a prefix minimum --- one running variable. The order of the two updates matters: \emph{first} use the minimum of $[0, j)$, \emph{then} add $a_j$ to it, otherwise $a_j$ could pair with itself.

\begin{diagrambox}
\begin{tikzpicture}
    \dpindex{pi}{1.35}{$j$}{0}{6}
    \dparray{p}{0.6}{\texttt{a}}{7/acell, 2/acell, 5/acell, 8/acell, 1/akeep, 6/aread, 3/acell}
    \aptr{p4}{lo}{successgreen}
    \aptr{p5}{j}{accentorange}
    \node[note, anchor=west, text width=4.4cm, align=left] at (6.7,0.6)
        {at $j = 5$ the minimum of $[0, 5)$ is $1$ (at index 4), so the pair gives $6 - 1 = 5$; the overall best is $8 - 2 = 6$, found at $j = 3$};
\end{tikzpicture}
\figcap{The prefix minimum is the best partner of every later element. When $j = 5$ is reached, $\mathit{lo}$ has dropped from $2$ to $1$; the best pair so far ($8 - 2$, found at $j = 3$) is kept in \textit{best}.}
\end{diagrambox}

%------------------------------------------------------------
\subsubsection*{Variations map}

\begin{diagrambox}
\begin{tikzpicture}[
    var/.style={rectangle, rounded corners=2pt, draw=patternborder, fill=patternbg,
                font=\scriptsize, align=center, text width=3.25cm, minimum height=1.2cm, inner sep=3pt}]
    \node[var, fill=answerfill, draw=accentorange, text width=5cm, font=\scriptsize\bfseries] (root) at (0,0)
        {Single-pass aggregate\\state $\gets$ \textsc{update}(state, $a_i$)};
    \node[var] (v1) at (-5.4,-2) {\textbf{top-$k$, small $k$}\\distinct or not (Problem~\ref{ssec:ch1-third-maximum}, E1.5)};
    \node[var] (v2) at (-1.8,-2) {\textbf{global, then per position}\\max of the others (Problem~\ref{ssec:ch1-advantage})};
    \node[var] (v3) at (1.8,-2)  {\textbf{best pair $i < j$}\\prefix extreme as partner (E1.3)};
    \node[var] (v4) at (5.4,-2)  {\textbf{streaks}\\current run, best run (E1.1, E1.6)};
    \node[var] (v5) at (-3.6,-3.6) {\textbf{positions of extremes}\\first/last index of min and max (E1.7)};
    \node[var] (v6) at (0,-3.6)    {\textbf{sum and extreme together}\\e.g.\ total $- n\cdot\min$ (E1.2)};
    \node[var] (v7) at (3.6,-3.6)  {\textbf{adjacent pairs}\\aggregate over $\max(a_i, a_{i+1})$ (E1.4)};
    \node[draw=patternborder, dashed, rounded corners=4pt, fit=(v1)(v4)(v5)(v7), inner sep=5pt] (grp) {};
    \draw[flow] (root.south) -- (grp.north);
\end{tikzpicture}
\figcap{Every box is one scan (or two) with $\bigO{1}$ state. The labels point to the solved problems and exercises of this pattern.}
\end{diagrambox}

\textbf{Twists and how to handle them}
\begin{itemize}[leftmargin=*]
    \item \textbf{Fewer than $k$ distinct values:} the statement decides the fallback (Problem~\ref{ssec:ch1-third-maximum} returns the maximum). Keep a counter of filled slots or a sentinel.
    \item \textbf{Ties for the extreme:} \texttt{>} keeps the first occurrence, \texttt{>=} the last. Decide which index the problem wants.
    \item \textbf{Negative values:} never start a maximum at $0$; for products, the two \emph{smallest} values (most negative) matter too (Exercise E1.5).
    \item \textbf{Huge sums:} $n \cdot \max|a_i|$ beyond $2^{31}$ needs \texttt{long long} (Lesson~\ref{ssec:ch1-cost}).
    \item \textbf{``Except itself'':} keep the two largest \emph{with duplicates}; the answer for $i$ is $m_2$ if $a_i = m_1$, else $m_1$.
\end{itemize}

%------------------------------------------------------------
\subsubsection*{Code templates}

The first template computes the common aggregates of the table in one pass; the second keeps the $k$ largest distinct values for a small $k$.

\begin{lstlisting}[style=cpp]
#include <bits/stdc++.h>
using namespace std;
typedef long long ll;

// Input: n, then a[0..n-1] (n >= 1). One pass computes several aggregates.
int main() {
    int n;
    if (!(cin >> n)) return 0;
    vector<ll> a(n);
    for (auto& x : a) cin >> x;

    ll sum = 0;
    ll mn = LLONG_MAX, mx = LLONG_MIN;
    int firstMax = -1, lastMax = -1;
    ll top1 = LLONG_MIN, top2 = LLONG_MIN;          // two largest, duplicates allowed
    ll minSoFar = LLONG_MAX, bestDiff = LLONG_MIN;   // max a[j] - a[i], i < j
    int run = 0, bestRun = 0;                        // longest block of equal values
    for (int i = 0; i < n; i++) {
        ll x = a[i];
        sum += x;
        mn = min(mn, x);
        if (x > mx) { mx = x; firstMax = i; }        // strict: keeps the FIRST maximum
        if (x == mx) lastMax = i;                    // every tie moves it: the LAST maximum
        if (x > top1) { top2 = top1; top1 = x; }     // shift down
        else if (x > top2) top2 = x;
        if (i > 0) bestDiff = max(bestDiff, x - minSoFar);   // use [0, i) first ...
        minSoFar = min(minSoFar, x);                         // ... then add a[i]
        run = (i > 0 && a[i] == a[i - 1]) ? run + 1 : 1;
        bestRun = max(bestRun, run);
    }
    cout << "sum " << sum << ", min " << mn << ", max " << mx
         << " (first at " << firstMax << ", last at " << lastMax << ")\n";
    cout << "two largest " << top1 << " " << top2 << ", best a[j]-a[i] " << bestDiff
         << ", longest run " << bestRun << "\n";
}
\end{lstlisting}

\begin{lstlisting}[style=cpp]
#include <bits/stdc++.h>
using namespace std;

// The k largest DISTINCT values seen so far, k small (2, 3, 4 ...): O(k) per element.
struct TopDistinct {
    int k;
    vector<long long> v;                       // sorted decreasingly, no duplicates, size <= k
    explicit TopDistinct(int k_) : k(k_) {}
    void add(long long x) {
        for (long long y : v) if (y == x) return;          // already present
        int p = 0;
        while (p < (int)v.size() && v[p] > x) p++;          // first position with v[p] < x
        if (p == k) return;                                 // smaller than all k kept values
        v.insert(v.begin() + p, x);
        if ((int)v.size() > k) v.pop_back();                // forget the (k+1)-th
    }
};

int main() {
    TopDistinct t3(3);
    int a[] = {4, 9, 4, 7, 1, 9, 6};
    for (int x : a) {
        t3.add(x);
        cout << "after " << x << ":";
        for (long long y : t3.v) cout << " " << y;
        cout << "\n";
    }
}
\end{lstlisting}

\begin{insightbox}[Check the templates by hand]
On input \texttt{9 / 4 9 2 2 2 9 -1 7 3} the first template prints sum $37$, min $-1$, max $9$ (first at index 1, last at index 5), two largest $9\ 9$, best $a_j - a_i = 7 - (-1) = 8$ and longest run $3$ (the three 2's). The second prints the states $(4)$, $(9,4)$, $(9,4)$, $(9,7,4)$, $(9,7,4)$, $(9,7,4)$, $(9,7,6)$ --- the same states as the trace of Problem~\ref{ssec:ch1-third-maximum}.
\end{insightbox}

%------------------------------------------------------------
\subsubsection*{Common mistakes}

\begin{warningbox}[Mistake 1 --- a maximum initialised to 0]
\begin{lstlisting}[style=snippet]
int mx = 0;      for (int x : a) mx = max(mx, x);   // WRONG for all-negative arrays: returns 0
int mx = a[0];   for (int x : a) mx = max(mx, x);   // RIGHT (n >= 1)
\end{lstlisting}
\end{warningbox}

\begin{warningbox}[Mistake 2 --- forgetting to shift]
\begin{lstlisting}[style=snippet]
if (x > m1) m1 = x;                   // WRONG: the old maximum is lost, not demoted to m2
if (x > m1) { m2 = m1; m1 = x; }      // RIGHT
else if (x > m2) m2 = x;
\end{lstlisting}
\end{warningbox}

\begin{warningbox}[Mistake 3 --- a sentinel that can occur in the data]
\texttt{INT\_MIN} as ``unset'' fails on the input $[1, 2, -2147483648]$: the third maximum really is $-2^{31}$, but the code believes the third slot is empty. Use a wider type, a flag, or start from real elements.
\end{warningbox}

\begin{warningbox}[Mistake 4 --- updating the state before using it]
For pairs $i < j$, updating the prefix minimum with $a_j$ \emph{before} computing $a_j - \min$ pairs $a_j$ with itself and reports $0$ for a strictly decreasing array where the answer should be ``no pair'' (or a negative value, depending on the statement).
\end{warningbox}

\begin{warningbox}[Mistake 5 --- distinct versus not]
Using the distinct-values update for ``the largest element other than $a_i$'' gives a wrong answer on $(4, 4, 4, 4)$: the strongest other participant has strength $4$, not ``none''. Decide from the statement whether equal values count separately.
\end{warningbox}

\begin{checklistbox}[``I'm stuck'' checklist for Pattern 1]
\begin{itemize}[leftmargin=*]
    \item[$\square$] Is the answer one number (or one global number used at every position)? Then try a scan.
    \item[$\square$] What exactly is the state? Write it as ``after reading $[0, i)$, the variables hold \ldots''.
    \item[$\square$] Can I update it from the old state and $a_i$ alone? If not, it is a different pattern.
    \item[$\square$] What is the state of the empty prefix? Can a sentinel collide with real data?
    \item[$\square$] Ties: distinct values or elements? First index or last?
    \item[$\square$] Pairs: am I using the prefix \emph{before} adding the current element?
    \item[$\square$] Could the sums overflow \texttt{int}?
\end{itemize}
\end{checklistbox}
